\documentstyle[% 


righttag% 


,11pt% 


,amssymb% 


,russian%               remove in Eng. 


,izvru%                 Set izven% in Eng. 


]{amsart} 


 


%  split line 685 with \notag


 


\numberwithin{equation}{section} 


\newcommand{\tts}[1]{} 


 


%\hoffset=-15mm%                  For Eng. 


\setcounter{page}{36} 


\god{2002} 


\nomer{4~(479)} %                        Remove in Eng. 


%\nomer{Vol.\,46, No.\,4}%               Set in Eng. 


%\str{\pageref{begin}--\pageref{end}}%  Set in Eng. 


\udk{517.933, 517.972} 


 


\author{\it Р.Г.\,Мухарлямов} 


% NB:   ^^ remove in Eng. 


\title{О численном решении уравнений экстремалей 


вариационной задачи с ограничениями} 


\thanks{Работа выполнена при финансовой поддержке Российского фонда 


фундаментальных исследований, код \No\,99-01-0164, и НТП Министерства


образования Российской Федерации, код \No\,205.02.01.038.} 


 


\date{} 


%\pagestyle{myheadings}%         Set in Eng. 


\pagestyle{plain} %              Remove in Eng. 


\begin{document} 


%\thispagestyle{plain}%          Set in Eng. 


\maketitle 


\label{begin} 


 


 


\section{Введение} 


 


Известная вариационная задача с ограничениями [1] состоит в определении 


экстремалей функционала 


\begin{gather} 


I=\int^{t_1}_{t_0}L(q,v,t)d t,\ \ q=(q_1,\dotsc,q_n), 


\ \ v=(v_1,\dotsc,v_n), \ \ v_i=\Dot q_i=\frac{d q_i}{d t}, 


\ \ i=1,\dotsc,n,\\ 


% 


q(t_0)=q^0,\quad q(t_1)=q^1, \\ 


% 


f(q,t)=0,\ \ f'(q,v,t)=0,\ \ f=(f_1,\dotsc,f_m), 


\ \ f'=(f_{m+1},\dotsc,f_r),\ \ r\le n, 


\end{gather} 


и сводится к определению вектор-функции $q=q(t)$, удовлетворяющей 


одновременно 


граничным условиям (1.2), уравнениям связей (1.3) и уравнениям 


Эйлера--Лагранжа 


\begin{gather} 


\Dot q=v,\quad \frac{d}{d t}\frac{\partial L}{\partial v} 


-\frac{\partial L}{\partial q}=F^T\lambda,\\ 


% 


\frac{\partial L}{\partial q}=\bigg( 


\frac{\partial L}{\partial q_1},\dotsc, 


\frac{\partial L}{\partial q_n}\bigg),\ \ F=(f_{\rho i}), 


\ \ F^T=(f_{i\rho}),\ \ f_{\mu i}= 


\frac{\partial f_\mu}{\partial q_i}, 


\ \ f_{\sigma i}=\frac{\partial f_\sigma}{\partial v_i},\notag \\ 


\rho=1,\dotsc,r,\ \ \mu=1,\dotsc,m,\ \ \sigma=m+1,\dotsc,r. \notag


\end{gather} 


 


Система (1.4) содержит вектор множителей Лагранжа 


$\lambda=(\lambda_1,\dotsc,\lambda_r)$, 


определяемый из условий $\Ddot f=0$, $\Dot f'=0$. При этом многообразие 


$\Omega(t)$ в пространстве 


переменных $(q,v)$, описываемое уравнениями 


\begin{gather} 


f(q,t)=0,\ \ \Dot f(q,v,t)=0,\ \ f'(q,v,t)=0,\\ 


\Dot f(q,v,t)\equiv f_q v+f_t,\ \ f_q=(f_{\mu i}), 


f_t=(f_{\mu t}),\ \ \mu=1,\dotsc,m,\ \ i=1,\dotsc,n, \notag


\end{gather} 


является устойчивым интегральным многообразием системы (1.3), 


(1.4). Однако оно не будет устойчивым асимптотически, и при 


численном решении системы уравнений (1.3), (1.4) отклонения от 


многообразия $\Omega(t)$ будут возрастать вследствие ошибок численного 


интегрирования. 


 


Уравнения (1.3), (1.4) составляют систему 


дифференциально-алгебраических уравнений индекса три (ДАУ-3). 


Индекс системы ДАУ определяется [2] числом, превышающим на 


единицу порядок максимальной производной от уравнений связей 


(1.3), необходимый для исключения из уравнений Эйлера--Лагранжа 


вектора $\lambda$. ДАУ являются в последнее время объектом интенсивных 


исследований [2], [3]. Основной проблемой в этих исследованиях 


является обеспечение асимптотической устойчивости интегрального 


многообразия $\Omega(t)$. Условия устойчивости интегрального многообразия 


ДАУ-2 были сформулированы в [4]. Возможные отклонения от 


уравнений связей (1.3), возникающие при численном решении 


системы (1.3), (1.4), можно ограничить. С этой целью в п.\,2 


вводятся необходимые определения и предлагается метод 


составления уравнений Эйлера--Лагранжа, позволяющий ограничить 


отклонения от уравнений связей при численном решении системы, 


состоящей из уравнений экстремалей и уравнений связей. В п.\,3 


строится алгоритм приведения системы ДАУ-3 к системе ДАУ-2. В п.\,4 


приводятся условия асимптотической устойчивости многообразия 


$\Omega(t)$ и в п.\,5 определяются условия, обеспечивающие устойчивость 


его при численном решении ДАУ-2 методом Эйлера и методом 


Рун\-ге-Кут\-та. 


 


\section{Уравнения экстремалей} 


 


При решении системы (1.3), (1.4) уравнения связей могут оказаться 


нарушенными уже в начальный момент. Если вектор $\lambda$ считать 


определенным из 


условий $\Ddot f=0$, $\Dot f'=0$, и систему дифференциальных уравнений 


(1.4) с начальными 


условиями $q(t_0)=q^0$, $v(t_0)=v^0$ решать численно, то вполне может 


оказаться, что при 


$t=t_0$ условия (1.5) будут нарушены: 


$$ 


f(q^0,t_0)=f^0,\ \ \Dot f(q^0,v^0,t_0)=\Dot f^0, 


\ \ f'(q^0,v^0,t_0)=f'{}^0. 


$$ 


Если даже начальные данные точно удовлетворяют уравнениям (1.5), 


то дальнейшее накопление ошибок численного интегрирования 


неизбежно приведет к росту отклонений от уравнений связей. 


Введением соответствующей коррекции в разностную схему решения 


системы (1.3), (1.4) отклонения от уравнений связей (1.3) можно 


ограничить. Для оценки отклонений от уравнений (1.5) запишем 


уравнения связей в виде равенств 


\begin{equation} 


f_\mu(q,t)=\alpha_\mu,\ \ \Dot f_\mu(q,v,t)=\alpha_{m+\mu}, 


\ \ f_\sigma(q,v,t)=\alpha_{m+\sigma}. 


\end{equation} 


Правые части уравнений (2.1) выбираются как решения 


$\alpha_\gamma=\alpha_\gamma(t)$ 


дифференциальных уравнений 


\begin{equation} 


\Dot\alpha_\gamma=\varphi_\gamma(\alpha,q,v,t), 


\ \ \varphi_\mu=\alpha_{m+\mu}, 


\ \ \alpha=(\alpha_1,\dotsc,\alpha_{m+r}), 


\ \ \varphi_\gamma(0,q,v,t)=0,\ \ \gamma=1,\dotsc,m+r, 


\end{equation} 


рассматриваемых совместно с уравнениями экстремалей (1.4) и 


удовлетворяющих 


начальным условиям 


\begin{gather} 


q_i(t_0)=q^0_i,\ \ v_i(t_0)=v^0_i, 


\ \ \alpha_\mu(t_0)=f^0_\mu,\ \ \alpha_{m+\mu}(t_0)=\Dot f^0_\mu, 


\ \ \alpha_{m+\sigma}(t_0)=f^0_\sigma,\\ 


% 


f^0_\mu=f_\mu(q^0,t_0),\ \ \Dot f^0_\mu=\sum^n_{i=1} 


f_{\mu i}(q^0,t_0)v^0_i+f_{\mu t}(q^0,t_0), 


\ \ f^0_\sigma=f_\sigma(q^0,v^0,t_0). \notag


\end{gather} 


 


Правые части системы (2.2) можно выбрать так, чтобы ее 


тривиальное решение $\alpha_1=\dotsb=\alpha_{r+m}=0$ было асимптотически 


устойчиво [4]. Но 


даже при этом отклонения от уравнений связей при численном 


интегрировании уравнений (1.4), (2.2) могут возрастать [2]. 


Стабилизацию уравнений связей (1.3) при численном интегрировании 


можно обеспечить за счет ограничений, накладываемых на правые 


части уравнений системы (2.2). Далее будет показано, что 


использование уравнений связей (2.1), (2.2) не изменяет 


структуры уравнений (1.4), но приводит к изменению выражений 


вектора $\lambda$, составленного из множителей Лагранжа. 


 


\begin{defn} 


Условия (2.1), наложенные на функции $q_1(t),\dotsc,q_n(t)$ и их 


производные $v_1(t),\dotsc,v_n(t)$, называются {\it уравнениями 


программных связей}, а условия 


(2.2), которым удовлетворяют их правые части 


$\alpha_1(t),\dotsc,\alpha_{m+r}(t)$, --- 


{\it уравнениями возмущений связей}. 


\end{defn} 


 


\begin{defn} 


Интегральное многообразие $\Omega(t)$ системы (1.4), описываемое 


уравнениями (1.5), называется {\it устойчивым интегральным 


многообразием}, если для 


любого $\varepsilon$ существует такое $\delta$, что для всех начальных 


значений (2.3), 


удовлетворяющих условию $\|\alpha^0\|\le\delta$, при всех $t>t_0$ будет 


выполняться неравенство 


$\|\alpha(t)\|\le\varepsilon$, и {\it асимптотически устойчивым 


интегральным многообразием}, 


если оно устойчиво и $\lim\limits_{t\to\infty}\|\alpha(t)\|=0$. 


\end{defn} 


 


\begin{thm} 


Для того чтобы функции $q_i=q_i(t)$, $i=1,\dotsc,n$, определяли 


экстремали 


функционала $(1.1)$ при ограничениях $(1.3)$, $(2.3)$, $f^0_\mu=\Dot 


f^0_\mu=0$, $\mu=1,\dotsc,m$, $f^0_\sigma=0$, $\sigma=m+1,\dotsc,r$, 


необходимо, чтобы они удовлетворяли системе уравнений $(1.4)$, $(2.1)$, 


$(2.2)$. 


\end{thm} 


 


{\bf Доказательство.} Условие стационарности 


$$ 


\delta\int^{t_1}_{t_0}(L+\lambda^T\wh f)d t=0,\quad 


\wh f=(f_1,\dotsc,f_r), 


$$ 


функционала (1.1) на экстремалях, удовлетворяющих условиям 


(2.1), (2.2), приводит к известному равенству 


\begin{equation} 


\int^{t_1}_{t_0}(\delta L+\lambda^T\delta\wh f)d t=0. 


\end{equation} 


 


Выражение вектора $\lambda$ определяется способом варьирования функций 


$f_\mu$ и $f_\sigma$. 


Положим [5] 


\begin{equation} 


\delta f_\mu=\sum^n_{i=1}\frac{\partial f_\mu}{\partial q_i} 


\delta q_i,\quad \delta f_\sigma=\sum^n_{i=1} 


\frac{\partial f_\sigma}{\partial v_i}\delta q_i. 


\end{equation} 


Тогда, учитывая условия (2.5), равенства $\delta q(t_0)=\delta q(t_1)=0$, 


получим


$$ 


\delta L=E^T(L)\delta q,\quad E(L)=\frac{d}{d t} 


\frac{\partial L}{\partial v}-\frac{\partial L}{\partial q},


$$ 


и принимая во внимание обозначения (1.4), представим равенство (2.4) в 


виде 


\begin{equation} 


\int^{t_1}_{t_0}(-E^T(L)+\lambda^T F)\delta q\,d t=0. 


\end{equation} 


 


Условие (2.6) выполняется, если 


\begin{equation} 


(E^T(L)-\lambda^T F)\delta q=0. 


\end{equation} 


Из (2.1), (2.5) следует равенство 


\begin{equation} 


F\delta q=\delta\wh\alpha,\quad \wh\alpha= 


(\alpha_1,\dotsc,\alpha_m,\alpha_{2m+1},\dotsc,\alpha_{m+r}), 


\end{equation} 


которое представляет собой систему линейных алгебраических уравнений 


относительно составляющих вектора $\delta q$ с прямоугольной матрицей 


коэффициентов. В 


[6] доказано, что общее решение уравнения (2.8) определяется выражением 


\begin{equation} 


\delta q=[FC]\delta s+F^+\delta\wh\alpha, 


\end{equation} 


где $[FC]$ есть векторное произведение векторов $f_{1q},\dotsc,f_{m 


q},f_{m+1v},\dotsc,f_{r v}$, составляющих 


строки матрицы $F$ и произвольных векторов $c_{r+1},\dotsc,c_{n-1}$, 


составляющих матрицу 


$C=(c_{\beta i})$, $\beta=r+1,\dotsc,n-1$, $F^+=F^T(F F^T)^{-1}$, 


$\delta s$ --- произвольная малая скалярная 


величина. Составляющая $a_{1k}$ векторного произведения 


$a_1=[a_2a_3\dotsc a_n]$ вычисляется как 


значение определителя, первая строка которого состоит из нулей, 


кроме единицы на $k$-ом месте, остальные строки составляют 


компоненты векторов $a_2,a_3,\dotsc,a_n$. 


 


Подставляя значение (2.9) в (2.7) и учитывая тождества $F[FC]\equiv0$, 


$F F^+\equiv I_r$, где 


$I_r$ --- единичная матрица порядка $r$, получим 


\begin{equation} 


E^T(L)[FC]\delta s+(E^T(L)F^+-\lambda^T I_r)\delta\wh\alpha=0. 


\end{equation} 


 


Так как скаляр $\delta s$ и вектор $\delta\wh\alpha$ произвольны, то из 


(2.10) 


следуют равенства 


\begin{equation} 


E(L)=F^T k,\quad k=\lambda, 


\end{equation} 


эквивалентные системе уравнений (1.4). 


 


Таким образом, если функции $q_i=q_i(t)$, $i=1,\dotsc,n$, удовлетворяющие 


начальным 


условиям (2.3), являются экстремалями функционала (1.1), то они 


должны представлять собой соответствующее решение системы 


уравнений (2.11) с вектором $\lambda$, определенным из равенств (2.1), 


(2.2). 


 


\section{Сведение ДАУ-3 к системе ДАУ-2} 


 


Система уравнений (2.1)--(2.3), (2.11) представляет собой 


систему ДАУ-3. При соответствующем выборе вектора $\lambda$ заменой 


переменных ДАУ-3 сводится к системе ДАУ-2, которая приводится к 


системе дифференциальных уравнений первого порядка, допускающей 


частные интегралы, определяемые уравнениями связей (1.5). 


 


\begin{thm} 


Система $(2.11)$ с вектором $\lambda=\lambda(q,v,t)$, определенным с 


учетом 


$(2.1)$, $(2.2)$, может быть представлена системой дифференциальных 


уравнений первого порядка 


\begin{equation} 


\Dot x=v^\tau(x,t)+J(x,t)g(x,t), 


\end{equation} 


допускающей частные интегралы 


\begin{equation} 


g(x,t)=0,\quad g=(f_1,\dotsc,f_m,\Dot f_1,\dotsc, 


\Dot f_m,f_{m+1},\dotsc,f_r). 


\end{equation} 


\end{thm} 


 


{\bf Доказательство.} Вычислив выражение левой части уравнения (2.11), 


представим 


систему (1.4) в виде 


\begin{gather} 


\frac{d q}{d t}=v,\quad \frac{d v}{d t}=M^{-1}(F^T\lambda-\gamma),\\ 


M=(m_{j h}),\quad m_{j h}=\frac{\partial^2L} 


{\partial v_j\partial v_h},\quad j,h=1,\dotsc,n,\notag \\ 


\gamma=(\gamma_1,\dotsc,\gamma_n),\quad \gamma_h=\sum^n_{j=1} 


\frac{\partial^2L}{\partial v_h\partial q_j}v_j+ 


\frac{\partial^2L}{\partial v_h\partial t}- 


\frac{\partial L}{\partial q_h}.\notag 


\end{gather} 


 


В частности, если $L=\frac{1}{2}v^T M v-V(q)$, $m_{j h}=m_{j h}(q)$, 


$j,h=1,\dotsc,n$ [5], то 


$$ 


\gamma_h=\sum^n_{k,j=1}\gamma_{k j,h}v_k v_j- 


\frac{\partial V}{\partial q_h},\quad 


\gamma_{k j,h}=\frac{1}{2}(m_{h k,j}+m_{j h,k}-m_{k j,h}),\quad 


m_{jh,s}=\frac{\partial m_{j h}}{\partial q_s}. 


$$ 


 


Запишем систему уравнений (2.1)--(2.2) в виде 


\begin{gather} 


f(q,t)=y,\quad \Dot f(q,v,t)=\Dot y,\quad f'(q,v,t)=y',\\ 


% 


\frac{d y}{d t}=\Dot y,\quad \frac{d\Dot y}{d t}=A_{10}y+ 


A_{11}\Dot y+A_{12}y',\quad \frac{d y'}{d t}=A_{20}y+ 


A_{21}\Dot y+A_{22}y', \\ 


A_{\alpha\beta}=A_{\alpha\beta}(x,t),\ \ \alpha=1,2, 


\ \ \beta=0,1,2. \notag 


\end{gather} 


 


В частности, если матрицы $A_{\alpha\beta}$, $\alpha=1,2$, $\beta=0,1,2$, 


являются постоянными 


блочно-диагональными, то система (3.5) соответствует линейной комбинации 


уравнений связей и их производных, предложенной в [7]. 


 


Введем обозначения 


\begin{alignat}{3} 


x&=(q,v), &\quad w&=(v,-\gamma), &\quad g&=(f,\Dot f,f'),\notag \\ 


N&= 


\begin{pmatrix} 


I_n & 0_{n,n} \\ 0_{n,n} & M 


\end{pmatrix} 


, &\quad D&= 


\begin{pmatrix} 


0_{m,n} & f_q \\ 0_{r-m,n} & f'_v 


\end{pmatrix} 


, &\quad G&= 


\begin{pmatrix} 


f_q & f_v \\ \Dot f_q & \Dot f_v \\ f'_q & f'_v 


\end{pmatrix} 


, 


\end{alignat} 


где $I_n$ --- единичная матрица, $0_{m,n}$ --- $m\times n$-матрица, 


состоящая из нулей, и запишем 


систему (3.3)--(3.5) в виде управляемой системы 


\begin{equation} 


\Dot x=N^{-1}w+N^{-1}D^T\lambda 


\end{equation} 


с вектором управления $\lambda$, соответствующим уравнению программных 


связей 


\begin{equation} 


g(x,t)=\alpha,\quad \alpha=(y,\Dot y,y'), 


\end{equation} 


и уравнению возмущений связей 


\begin{equation} 


\Dot\alpha=\wt A\alpha,\quad \wt A= 


\begin{pmatrix} 


0_{m, m} & I_m & 0_{m,r-m} \\ A_{10} & A_{11} & A_{12} \\ 


A_{20} & A_{21} & A_{22} 


\end{pmatrix} 


. 


\end{equation} 


 


Дифференцирование выражения (3.8) с учетом обозначений (3.6) и уравнений 


(3.7), (3.9) приводит к равенству 


\begin{equation} 


G\Dot x+g_t=\wt A g, 


\end{equation} 


которое состоит из очевидного тождества $\Dot f-f_q v-f_t\equiv0$ 


и уравнения для определения вектора $\lambda$ множителей Лагранжа 


системы (3.3)--(3.5): 


\begin{gather} 


S\lambda=A g-s,\quad S=F M^{-1}F^T,\\ 


A= 


\begin{pmatrix} 


A_{10}A_{11}A_{12} \\ A_{20}A_{21}A_{22} 


\end{pmatrix} 


,\quad s= 


\begin{pmatrix} 


\Dot f_q v-f_q M^{-1}\gamma+\Dot f_t \\ 


f'_q v-f'_v M^{-1}\gamma+f'_t 


\end{pmatrix} 


.\notag 


\end{gather} 


 


Подставляя в (3.7) решение $\lambda=S^{-1}(A g-s)$ уравнения (3.11), 


получим уравнение 


(3.1), которое содержит в правой части вектор $\Dot x^r=N^{-1}(w-D^T 


S^{-1}s)$ и матрицу 


$J=N^{-1}D^T S^{-1}A$, удовлетворяющие равенствам 


\begin{equation} 


G \Dot x^r\equiv(\Dot f,0_{r1}),\quad G J=(0_{m,m+r},A). 


\end{equation} 


 


Из (3.10), (3.12) следует, что вектор $\Dot x^r$ с точностью до 


произвольного множителя $c$ может быть представлен векторным 


произведением $\Dot x^r=c[G C]$, $J=G^+ A$, т.\,е. уравнение (3.1) 


содержится во 


множестве систем дифференциальных уравнений [4] 


$$ 


\Dot x=c[G C]+G^+(A g-g_t),\quad 


G^+=G^T W^{-1},\quad W=G G^T, 


$$ 


и выполняется равенство $\Dot g=A g$. 


 


Если $f_t=0$, $g=(f,\Dot f)$, то, полагая 


$$ 


\Dot g=-\gamma L(Z L)^{-1}g,\quad Z= 


\begin{pmatrix} 


F & 0 \\ v^T F_{q^T} & F 


\end{pmatrix} 


,\quad L= 


\begin{pmatrix} 


F^T & 0 \\ 0 & F^T 


\end{pmatrix} 


,\quad F=f_q, 


$$ 


приходим к системе, исследованной в [2]. 


 


\section{Устойчивость связей} 


 


Условия асимптотической устойчивости интегрального многообразия 


$\Omega(t)$ системы дифференциальных уравнений (3.1), заданного 


уравнением (3.2), сформулированы в [4]. Проведенное 


преобразование ДАУ-3 (1.4), (2.1), (2.2) к ДАУ-2 позволяет 


использовать известные критерии для определения условий 


асимптотической устойчивости интегрального многообразия $\Omega(t)$ 


системы (1.4), (2.1), (2.2). Определим в момент времени $t$ 


расстояние $\rho(x,\Omega(t))$ от точки $x\in R_{2n}$ пространства 


переменных $(q_1,\dotsc,q_n,v_1,\dotsc,v_n)$ до 


интегрального многообразия $\Omega(t)\in R_{2n}$ равенством 


$\rho(x,\Omega(t))=\|\alpha(t)\|$, $\|\alpha\|=\sqrt{\alpha^T\alpha}$. 


Тогда из 


асимптотической или экспоненциальной устойчивости тривиального 


решения $\alpha_1=\dotsb=\alpha_{r+m}=0$ системы (2.2) следует 


соответственно асимптотическая 


или экспоненциальная устойчивость интегрального многообразия $\Omega(t)$ 


системы (2.11). 


 


Для системы (3.9) с постоянной матрицей $\wt A$ справедлива 


 


\begin{thm} 


Если $\rho(x,\Omega(t))=\|\alpha(t)\|$ и все корни характеристического 


уравнения 


$D(\varkappa)\equiv\det(\wt A-\varkappa I_{m+r})=0$ системы $(3.9)$ имеют 


отрицательные действительные части, то интегральное 


многообразие $\Omega(t)$ системы $(3.7)$ асимптотически устойчиво. 


\end{thm} 


 


Используя метод функций Ляпунова, сформулируем условия асимптотической 


устойчивости многообразия $\Omega(t)$. 


 


\begin{thm} 


Если $\rho(x,\Omega(t))=\|\alpha(t)\|$ и для уравнений возмущений связей 


$(3.9)$ 


существует знакоопределенная функция $V(\alpha,x,t)$, допускающая 


бесконечно малый 


высший предел, производная которой $\Dot V(\alpha,x,t)$, вычисленная в 


силу уравнений $(3.7)$, 


$(3.9)$, является знакоопределенной функцией противоположного знака, то 


интегральное многообразие 


$\Omega(t)$ системы $(3.7)$ асимптотически устойчиво. 


\end{thm} 


 


В частности, если $2V=\alpha^T L\alpha$, $L=L^T=(l_{i j})$, 


$i,j=1,\dotsc,m+r$, --- положительно 


определенная квадратичная форма с постоянными коэффициентами, то $\Dot 


V=\alpha^T L\wt A\alpha$, и 


может быть использован критерий Сильвестра положительной определенности 


квадратичной формы $-\Dot V=\alpha^T A'\alpha$, $A'=-L\wt A=(a'_{i j})$: 


$$ 


a'_{11}(x,t)\ge\varepsilon>0,\quad 


\begin{vmatrix} 


a'_{11}(x,t) & a'_{12}(x,t) \\ a'_{21}(x,t) & a'_{22}(x,t) 


\end{vmatrix} 


\ge\varepsilon>0,\dotsc, 


\begin{vmatrix} 


a'_{11}(x,t) & \dotsc & a'_{1,m+r}(x,t) \\ 


\hdotsfor[1.5]{3} \\ 


a'_{m+r,1}(x,t) & \dotsc & a'_{m+r,m+r}(x,t) 


\end{vmatrix} 


\ge\varepsilon>0. 


$$ 


 


Если на переменные $q$, $v$ наложены только связи вида $f'(q,v,t)=0$, то 


$m=0$, 


$\wt A=A_{22}(x,t)$ и оказывается применимой теорема об экспоненциальной 


устойчивости 


многообразия $\Omega(t)$ [8] системы (3.7). 


 


\begin{thm} 


Если 


\begin{itemize} 


\item[1)] $\rho(x,\Omega(t))=\|\alpha(t)\|$, $\alpha(t)=f'(x,t);$ 


\item[2)] $\Dot\alpha=-k(x,t)W(x,t)L\alpha$, скалярная функция $k(x,t)$ 


ограничена снизу$:$ $k(x,t)\ge k_0>0$, 


$W=f'_v(f'_v)^T$, матрица $L$ является постоянной и удовлетворяет условиям 


Сильвестра$;$ 


\item[3)] существует такое $w_0>0$, что при всех $x=x(t)\in\Omega_h(t)$, 


$t\ge t_0$, где $\Omega_h(t)$ --- 


$h$-окрестность многообразия $\Omega(t)$, выполняется неравенство 


$\alpha^T W(x,t)\alpha\ge w^2_0\|\alpha\|^2;$ 


\item[4)] $l_1$, $l_m$ --- наименьшее и наибольшее собственные значения 


матрицы $L$, 


$\gamma=g(x_0,t_0)\sqrt{l_m/l_1}$, $\beta=k_0w_0l^2_1/l_m$, 


\end{itemize} 


то интегральное многообразие 


$\Omega(t)$ системы $(3.7)$ экспоненциально устойчиво и 


выполняется условие $\sqrt{\alpha^T(t)\alpha(t)}\le\gamma 


e^{-\beta(t-t_0)}$. 


\end{thm} 


 


\section{Численное решение уравнений экстремалей} 


 


Уравнения (1.4), (2.1), (2.2), если даже обеспечена асимптотическая 


устойчивость многообразия $\Omega(t)$, не могут гарантировать при 


численном решении 


необходимую точность выполнения равенства (1.5). Для оценки 


точности численного решения уравнений экстремалей воспользуемся 


системой ДАУ-2, представленной уравнениями (3.1), (3.2). 


 


Пусть известны начальные значения $t_0$, $x^0$, для которых 


$\|g^0\|=\sup\limits_i |g^0_i|<\varepsilon$, $g(x^0,t_0)=g^0$ 


и построено уравнение (3.1), допускающее асимптотически устойчивое 


интегральное многообразие $\Omega(t)$, 


заданное уравнением (3.2). Полагая $a=A(x,t)g$ и 


используя правую часть системы (3.1), построим разностное уравнение 


\begin{equation} 


x^{k+1}=x^k+\tau v^k,\quad x^k=x(t_k),\quad t_{k+1}=t_k+\tau, 


\quad v=v^\tau(x,t)+J(x,t)g(x,t). 


\end{equation} 


 


\begin{thm} 


Существуют такие постоянные $\alpha$, $\tau_1$, $\varepsilon$ и матрица 


$A(x,t)$, что при выполнении неравенств 


 


$1)$ $\|g^0\|\le\varepsilon;$ 


 


$2)$ $\tau\le\tau_1;$ 


 


$3)$ $\|E+\tau A(x,t)\|\le\alpha<1;$ 


 


$4)$ $\frac{\tau^2_1}{2}\|g^{(2)}\|\le(1-\alpha)\varepsilon$, 


$g^{(2)}=v^T g_{x x}v+2g_{x t}+g_{t t}$, 


 


\noindent{}решение разностного уравнения $(5.1)$ будет удовлетворять 


условию 


\begin{equation} 


\|g^k\|\le\varepsilon\ \ \forall k=1,\dotsc,K. 


\end{equation} 


\end{thm} 


 


{\bf Доказательство.} Пусть для решения уравнения (3.1) используется 


разностная 


схема (5.1) и выполнены условия 1)--4). Предположим, что условие 


(5.2) выполняется при некотором значении $k$, и представим вектор 


$g^{k+1}=g(x^{k+1},t_{k+1})$ разложением в ряд 


\begin{gather} 


g^{k+1}=g^k+g^k_x\Delta x^k+g^k_t\tau+\frac{\tau^2}{2} 


g^{(k2)},\quad \Delta x^k=\tau v^k, \\ 


% 


g^{(k2)}=(g^{(k2)}_1,\dotsc,g^{(k2)}_{m+r}),\quad 


g^{(k2)}_\gamma=\frac{1}{2\tau^2}\bigg(\sum_{p,q} 


\wt g^{\,(k)}_{\gamma,p q}\Delta x^k_p \Delta x^k_q+2\tau\sum_p 


\wt g^{\,(k)}_{\gamma,p t}\Delta x^k_p+\tau^2 


\wt g^{\,(k)}_{\gamma,t t}\bigg),\notag 


\end{gather} 


$\wt g^{\,(k)}_{\gamma,p q}$, 


$\wt g^{\,(k)}_{\gamma,p t}$, 


$\wt g^{\,(k)}_{\gamma,t t}$ 


--- значения частных производных 


$\frac{\partial^2g_\gamma}{\partial x_p \partial x_q}$, 


$\frac{\partial^2g_\gamma}{\partial x_p \partial t}$, 


$\frac{\partial^2g_\gamma}{\partial t^2}$, 


вычисленные 


при $x=x^k+\Theta\Delta x^k$, $t=t_k+\theta\tau$. Матрица 


$\Theta=(\theta_{j l})$, $0\le\theta_{j l}\le1$, и скаляр $\theta$, 


$0\le\theta\le1$, 


принимают значения, соответствующие промежуточным значениям производных. 


 


Далее, учитывая равенство $\Dot g=A(x,t)g$, запишем выражение (5.3) в виде 


$$ 


g^{k+1}=(E+\tau A^k)g^k+\frac{\tau^2}{2}g^{(k2)}. 


$$ 


 


Оценивая правую часть выражения (5.3) с учетом 1)--4), получим 


$$ 


\|g^{k+1}\|\le\|E+\tau A^k\|\,\|g^k\|+\frac{1}{2}\tau^2 


\|g^{(k2)}\|\le\alpha\varepsilon+(1-\alpha) 


\varepsilon=\varepsilon. 


$$ 


 


Определим условия, которые следует наложить на правую часть уравнений 


(3.9) 


для обеспечения точности (5.2) при численном решении уравнения (3.1) 


методом 


Рун\-ге-Кут\-та. Будем аппроксимировать решение уравнения (3.1) разностной 


схемой 


второго порядка точности 


\begin{gather} 


x^{k+1}=x^k+\Delta x^k,\\ 


\Delta x^k=\tau(1-\sigma)v^k+\sigma\ov v^k,\ \ \sigma>0,\\ 


\ov v^k=v(\ov x^k,t_k+\alpha\tau),\ \ \alpha>0,\\ 


\ov x^k=x_k+\alpha\tau v^k,\quad 


v=v^\tau(x,t)+J(x,t)g(x,t), 


\end{gather} 


$\alpha$, $\sigma$ --- параметры разностной схемы. Будем считать, что 


начальные условия 


удовлетворяют неравенству $\|g^0\|\le\varepsilon$. 


 


Представим вектор $g^{k+1}$ в виде суммы по степеням $\Delta x^k$, $\tau$ 


\begin{gather} 


g^{k+1}=g^k+g^k_x\Delta x^k+g^k_t\tau+\frac{1}{2} 


((\Delta x^k)^T g^k_{x^T x} \Delta x^k+ 2g^k_{x t}\Delta x^k\tau+ 


g^k_{t t}\tau^2)+R^{(k3)}_g,\\ 


% 


R^{(k3)}_g=(R^{(k3)}_{g1},\dotsc,R^{(k3)}_{g,m+r}),\notag \\ 


% 


R^{(k3)}_{g\gamma}=\frac{1}{3!}\bigg(\sum_{p,q,r} 


\wt g^{\,(k)}_{\gamma,p q r} \Delta x^k_p \Delta x^k_q 


\Delta x^k_r+3\tau\sum_{p,q}\wt g^{\,(k)}_{\gamma,p q t} 


\Delta x^k_p \Delta x^k_q+3\tau^2\sum_p 


\wt g^{\,(k)}_{\gamma,p t t}\Delta x^k_p +\tau^3 


\wt g^{\,(k)}_{\gamma,t t t}\bigg),\notag 


\end{gather} 


$\wt g^{\,(k)}_{\gamma,p q r}$, 


$\wt g^{\,(k)}_{\gamma,p q t}$, 


$\wt g^{\,(k)}_{\gamma,p t t}$, 


$\wt g^{\,(k)}_{\gamma,t t t}$ 


--- значения частных производных $\frac{\partial^3g_\gamma}{\partial 


x_p\partial x_q\partial x_r}$, $\frac{\partial^3g_\gamma}{\partial 


x_p\partial x_q\partial t}$, $\frac{\partial^3g_\gamma}{\partial 


x_p\partial t^2}$, $\frac{\partial^3g_\gamma}{\partial t^3}$,


вычисленные при


$x=x^k+\Theta\Delta x^k$, $t=t_k+\theta\tau$. Из (5.5), (5.6) следует 


равенство $\Delta x^k=\tau((1-\sigma)v^k+\sigma v(x^k+\alpha\tau 


v^k,t_k+\alpha\tau))$, 


кoтopoe в результате разложения в ряд может быть представлено в виде 


\begin{gather} 


\Delta x^k=\tau(v^k+\alpha\sigma\tau\Dot v^k)+R^{(k3)}_v,\\ 


R^{(k3)}_v=(R^{(k3)}_{v1},\dotsc,R^{(k3)}_{v n}),\notag \\ 


R^{(k3)}_{v j}=\frac{1}{2!}\alpha^2\sigma\tau^2\bigg( 


\sum_{p,q}\wt v^{\,(k)}_{j,p q}v^k_q v^k_q+2\sum_p 


\wt v^{\,(k)}_{j,p t}v^k_p+\wt v^{\,(k)}_{j,t t}\bigg),\notag 


\end{gather} 


$\wt v^{\,(k)}_{j,p q}$, 


$\wt v^{\,(k)}_{j,p t}$, 


$\wt v^{\,(k)}_{j,t t}$ 


--- значения частных производных $\frac{\partial^2v_j}{\partial 


x_p\partial x_q}$, $\frac{\partial^2 v_j}{\partial x_p\partial t}$, 


$\frac{\partial^2v_j}{\partial t^2}$, вычисленные 


при $x=x^k+\Theta\Delta x^k$, $t=t_k+\theta\tau$. 


 


Последовательное вычисление правой части (5.8) с учетом (5.9) позволяет 


выразить $g^{k+1}$ через $g^k$, $x^k$, $t_k$, $A^k$, $\Dot A^k$. Для 


этого подставим в (5.8) выражение (5.9) 


\begin{multline} 


g^{k+1}=g^k+\tau(g^k_x(v^k+\alpha\sigma\tau\Dot v^k+ 


R^{(k3)}_v)+g^k_t)+ 


\frac{\tau^2}{2}((v^k+\alpha\sigma\tau\Dot v^k+R^{(k3)}_v)^T 


g^k_{x^T x}(v^k+\alpha\sigma\tau\Dot v^k+R^{(k3)}_v)+ \\


+2g^k_{x t}(v^k+\alpha\sigma\tau\Dot v^k+R^{(k3)}_v)+ 


g^k_{t t})+R^{(k3)}_g 


\end{multline} 


и представим (5.10) в виде 


\begin{gather} 


g^{k+1}=g^k+\tau(g^k_x v^k+g^k_t)+\alpha\sigma\tau^2 


(g_x\Dot v)^k+\frac{\tau^2}{2}((v^k)^T g^k_{x^T x}v^k+ 


2g^k_{x t}v^k+g^k_{t t})+R^{(k3)}, \\ 


% 


R^{(k3)}=\tau g^k_x R^{(k3)}_v+\frac{\tau^2}{2} 


(\alpha\sigma\tau\Dot v^k+R^{(k3)}_v)^T g^k_{x^T x}


(\alpha\sigma\tau\Dot v^k


+ R^{(k3)}_v)+ \notag \\


+\tau^2(v^{k^T}g^k_{x^T x}(\alpha\sigma\tau\Dot v^k+R^{(k3)}_v)+ 


g^k_{x t}(\alpha\sigma\tau\Dot v^k+R^{(k3)}_v))+R^{(k3)}_f.\notag 


\end{gather} 


 


Далее, учитывая равенства 


$$ 


\Dot g=A(x,t)g,\quad (g_x\Dot v)=\Ddot g-((v)^T g_{x^T x}v 


+2g_{x t}v+g_{t t}), 


$$ 


будем иметь из (5.11) выражение 


\begin{equation} 


g^{k+1}=(E+\tau A^k+\tau^2\alpha\sigma((A^k)^2+\Dot A^k))g^k+ 


\frac{\tau^2}{2}(1-2\alpha\sigma)((v^k)^T g^k_{x^T x}v^k+ 


2g^k_{x t}v^k+g^k_{t t})+R^{(k3)}. 


\end{equation} 


 


Выберем параметры $\alpha$, $\sigma$ так, чтобы $2\alpha\sigma=1$. 


Тогда из (5.12) следует 


\begin{equation} 


g^{k+1}=\bigg(E+\tau A^k+\frac{1}{2}\tau^2((A^k)^2+ 


\Dot A^k)\bigg)g^k+R^{(k3)}. 


\end{equation} 


 


\begin{thm} 


Если для решения уравнения $(3.1)$ используется разностная схема 


Рун\-ге-Кут\-та второго порядка точности $(5.4)$--$(5.7)$, и при всех 


$x=x_k$, $t=t_k$, $k=1,\dotsc,K$, 


величины $\tau$, $q>0$, $\alpha$, $\sigma$, $R^{(k3)}$, $g^0$, $A(x,t)$ 


удовлетворяют условиям 


 


$1)$ $\|g^0\|\le\varepsilon;$ 


 


$2)$ $2\alpha\sigma=1;$ 


 


$3)$ $\|R^{(k3)}\|\le(1-q)\varepsilon;$ 


 


$4)$ $\big\|E+\tau A^k+\frac{1}{2}\tau^2((A^k)^2+\Dot A^k)\big\|\le q<1$, 


 


\noindent{}то выполняется неравенство $(5.2)$ при всех $k=1,\dotsc,K$. 


\end{thm} 


 


{\bf Доказательство.} Если неравенство (5.2) верно при некотором $k$, 


то из оценки 


правой части выражения (5.13) следует, что будет верно и неравенство 


$$ 


\|g^{k+1}\|\le\bigg\|E+\tau A^k\frac{1}{2}\tau^2((A^k)^2+ 


\Dot A^k)\bigg\|\,\|g^k\|+\|R^{(k3)}\|\le q\varepsilon+ 


(1-q)\varepsilon=\varepsilon. 


$$ 


 


Напомним, что для вычислений по разностной схеме (5.4)--(5.7) обычно 


используются два варианта 


\begin{alignat*}{4} 


1)\ \; \sigma&=1, &\ \; \alpha&=\frac{1}{2}: &\ \; 


\ov x^k&=x^k+\frac{\tau}{2}v(x^k,t_k), &\ \; x^{k+1}&=x^k+\tau v 


\bigg(\ov x^k,t_k+\frac{\tau}{2} \bigg), \\ 


% 


2)\ \; \sigma&=\frac{1}{2} &\ \; \alpha&=1: &\ \; \ov x^k&= 


x^k+\tau v(x^k,t_k), &\ \; x^{k+1}&= x^k+\frac{\tau}{2} 


(v(x^k,t_k)+v(\ov x^k,t_k)). 


\end{alignat*} 


 


 


\begin{thebibliography}{9} 


 


\bibitem{1} 


Блисс Г.А. {\em Лекции по вариационному исчислению}. -- М.: Ин. лит., 


1950. -- 347 с. 


\bibitem{2} 


Ascher U.M., Hongsheng Chin, Petzold L.R., Reich S. 


{\em Stabilization of constrained mechanical systems with DAEs and 


invariant manifolds} // J. Mech. of Structures and Machines. -- 


1995. -- V.\,23. -- P.\,135--158. 


\bibitem{3} 


Rentrop P., Strehmel K., Weiner R. {\em Ein \"Uberblick \"uber 


Einschrittverfaren zur numerischen Integration in der 


technischen Simulation} // GAMM-Mitteilungen. -- 1996. -- 


Bd.\,19. -- H.\,1. -- S.\,9--43. 


\bibitem{4} 


Мухарлямов Р.Г. {\em О построении множества систем дифференциальных 


уравнений 


устойчивого движения по интегральному многообразию} // Дифференц. 


уравнения. --


1969. -- Т.\,5. -- \No\,4. -- С.\,688--699. 


\bibitem{5} 


Бухгольц Н.Н. {\em Основной курс теоретической механики}. Ч.\,2. -- М.: 


Наука, 


1966. -- 332 с. 


\bibitem{6} 


Мухарлямов Р.Г. {\em О построении дифференциальных уравнений оптимального 


движения по заданному многообразию} // Дифференц. уравнения. -- 


1971. -- Т.\,7. -- \No\,10. -- С.1825--1834. 


\bibitem{7} 


Baumgarte J. {\em Stabilization of constraints and integrals of 


motion in dynamical systems} // Comp. Math. Appl. Mech. Eng. -- 


1972. -- V.\,1. -- P.\,1--16. 


\bibitem{8} 


Мухарлямов Р.Г. {\em Численное моделирование в задачах механики} // 


Вестник 


РУДН. Cер. прикл. матем. и механ. -- 1995. -- \No\,1. -- C.\,13--28. 


 


\end{thebibliography} 


 


\vskip 2cm 


 


\post{\it Российский университет}{\it Поступила} 


\post{\it Дружбы народов}{26.08.1999} 


 


\label{end} 


\end{document} 


